
\documentclass[a4paper,11pt]{ltjsarticle}
\usepackage[haranoaji]{luatexja-preset}
\usepackage{geometry}
\geometry{top=24mm,bottom=24mm,left=24mm,right=24mm}
\usepackage{setspace}
\setstretch{1.18}
\usepackage{fancyhdr}
\usepackage{longtable}
\usepackage{booktabs}
\usepackage{array}
\usepackage{tabularx}
\usepackage{enumitem}
\usepackage{tcolorbox}
\tcbuselibrary{breakable,skins}
\usepackage{hyperref}
\hypersetup{hidelinks,unicode=true,pdftitle={第02章 ソフトウェアアーキテクチャ 日本語まとめ},pdfauthor={OpenAI}}
\setlength{\parindent}{1em}
\setlength{\parskip}{0.25em}
\setlist[itemize]{leftmargin=2em,itemsep=0.2em,topsep=0.2em}
\setlist[enumerate]{leftmargin=2em,itemsep=0.2em,topsep=0.2em}
\newcolumntype{Y}{>{\raggedright\arraybackslash}X}
\newcolumntype{P}[1]{>{\raggedright\arraybackslash}p{#1}}
\pagestyle{fancy}
\fancyhf{}
\fancyhead[L]{第 02 章 ソフトウェアアーキテクチャ}
\fancyhead[R]{SWEBOK v4.0a 日本語まとめ}
\fancyfoot[C]{\thepage/\pageref{LastPage}}
\renewcommand{\headrulewidth}{0.4pt}
\usepackage{lastpage}

\newtcolorbox{statementbox}[1]{enhanced,breakable,colback=white,colframe=black,boxrule=0.6pt,arc=0mm,left=1mm,right=1mm,top=1mm,bottom=1mm,title=\textbf{#1},fonttitle=\normalsize,coltitle=black,colbacktitle=white,attach boxed title to top left={xshift=1mm,yshift=-1mm},boxed title style={colback=white,colframe=white,arc=0mm,boxrule=0pt}}
\newcommand{\todofig}[2]{\begin{statementbox}{TODO（図） #1}\textbf{画像生成用プロンプト}\par #2\end{statementbox}}
\newcommand{\keyterm}[1]{\textbf{#1}}
\newcommand{\en}[1]{（#1）}

\begin{document}
\thispagestyle{plain}
\begin{center}
{\Large\bfseries 第 02 章 ソフトウェアアーキテクチャ}\par\vspace{1em}
{\large Software Architecture の日本語まとめ}\par\vspace{0.7em}
{SWEBOK Guide v4.0a Chapter 02 を初学者向けに再構成}\par\vspace{1.2em}
{作成日：2026 年 4 月 29 日}\par
\end{center}
\vspace{1em}
\hrule
\vspace{0.7em}

\noindent\textbf{この資料の主張}\par
ソフトウェアアーキテクチャは、単なる「大まかな設計図」ではない。アーキテクチャは、ソフトウェア集約システムにおいて根本的であり、後から変えにくく、多くの品質や費用に影響する構造・関係・原則・重要な意思決定の集合である。よいアーキテクチャは、利害関係者の関心事を整理し、設計・実装・テスト・運用・保守に関わる人々が同じ理解を持つための土台になる。

\begin{statementbox}{この資料の読み方}
専門用語は原則として日本語で示し、必要に応じて英語を括弧内に併記する。英語名は、元資料、規格、ツール、文献で同じ概念を確認するための手がかりとして残す。初学者が前提知識なしに読めるように、各節では「何のための概念か」「実務では何を判断するのか」を重視する。
\end{statementbox}

\tableofcontents
\clearpage

\section*{全体概要：この章で学ぶこと}
\addcontentsline{toc}{section}{全体概要：この章で学ぶこと}
この章は、ソフトウェアの根本的な形を決め、関係者が共有し、後工程の判断に使うための「アーキテクチャ」を扱う。ここでいうアーキテクチャは、クラス図や構成図そのものではない。図や文書は、アーキテクチャを表すための手段である。アーキテクチャの中心にあるのは、主要な構成要素、それらの関係、設計と進化の原則、品質特性に関するトレードオフ、そして重要な意思決定である。

\noindent\textbf{到達目標}\quad この章を読み終えると、次のことができるようになる。
\begin{itemize}
  \item 「アーキテクチャ」という語が、分野、活動、成果物の三つの意味で使われることを説明できる。
  \item 利害関係者と関心事を分けて考え、アーキテクチャが誰の何に答えるものかを説明できる。
  \item ビューとビューポイントの違いを説明し、複数の図やモデルを使う理由を説明できる。
  \item アーキテクチャ様式、アーキテクチャパターン、参照アーキテクチャを区別できる。
  \item アーキテクチャ設計を、分析、統合、評価の反復活動として説明できる。
  \item アーキテクチャ評価、レビュー、測度の目的を説明できる。
\end{itemize}

\noindent\textbf{学習順序}\quad 初学者は、まず「アーキテクチャとは何か」と「利害関係者と関心事」を押さえる。その後、「どう記述するか」「どう設計するか」「どう評価するか」という流れで読むとよい。

\todofig{02 章の全体地図を入れる。}{白背景、モノクロ、A4 本文幅に収まる横長の学習用図解を作成する。中央に「ソフトウェアアーキテクチャ」を置き、周囲に「基礎」「アーキテクチャ記述」「アーキテクチャ過程」「アーキテクチャ評価」を配置する。各領域の下に「関心事」「ビュー」「分析・統合・評価」「レビュー・測度」などの小項目を置く。ラベルは日本語、細線、講義ノート風、余白を広めにする。}

\subsection*{最初に必要な用語}
\addcontentsline{toc}{subsection}{最初に必要な用語}
\begin{longtable}{P{0.28\linewidth}P{0.64\linewidth}}
\toprule
用語 & 初学者向けの説明 \\
\midrule
ソフトウェア集約システム & ソフトウェアが主要な役割を持つシステムである。ソフトウェアだけでなく、ハードウェア、人、手順、外部サービスも含み得る。 \\
アーキテクチャ & システムの根本的な概念、構成要素、関係、設計と進化の原則である。単なる図ではなく、重要な判断の集合でもある。 \\
利害関係者 & システムに影響を与える、または影響を受ける人や組織である。顧客、利用者、開発者、運用者、保守担当、監査者などが含まれる。 \\
関心事 & 利害関係者が気にする論点である。費用、納期、性能、信頼性、安全性、セキュリティ、使いやすさ、保守しやすさなどである。 \\
品質特性 & システムがどの程度よく振る舞うかを表す性質である。可用性、変更容易性、性能、信頼性、試験容易性などが例である。 \\
アーキテクチャ記述 & アーキテクチャを関係者が使えるように、文書、図、表、モデル、意思決定記録などで表した成果物である。 \\
ビュー & アーキテクチャの一面を表す図やモデルである。論理ビュー、配置ビュー、開発ビューなどがある。 \\
ビューポイント & ビューを作り、読み、使うための規約である。記法、表す要素、関心事、想定読者などを定める。 \\
アーキテクチャ上重要な要求 & アーキテクチャに影響する要求である。英語では Architecturally Significant Requirement、略して ASR と呼ぶ。 \\
アーキテクチャ技術的負債 & 今のアーキテクチャ判断が、将来の変更・運用・保守に追加費用を生む状態である。 \\
\bottomrule
\end{longtable}

\section*{導入：なぜソフトウェアアーキテクチャを独立して学ぶのか}
\addcontentsline{toc}{section}{導入：なぜソフトウェアアーキテクチャを独立して学ぶのか}
SWEBOK v4.0a では、ソフトウェアアーキテクチャをソフトウェア設計から独立した知識領域として扱う。これは、1990 年代以降にアーキテクチャ研究と実務が大きく発展し、設計の一部としてだけでは説明しきれない知識が増えたためである。

ソフトウェアアーキテクチャは、概念、表現、作業成果物、文脈、過程、方法、分析、評価という複数の視点から理解する必要がある。特に大規模なシステムでは、個々の実装詳細よりも先に、主要な構成要素、構成要素間の関係、外部環境との相互作用、品質特性への対応、変更に対する耐性を考える必要がある。

\begin{statementbox}{重要}
アーキテクチャは「実装の前に一度だけ描く図」ではない。要求、品質、組織、運用、保守、将来の進化に関わる継続的な判断の集合である。
\end{statementbox}

\section{ソフトウェアアーキテクチャの基礎}

\subsection{「アーキテクチャ」という語の意味}
ソフトウェア工学では、「アーキテクチャ」という語が少なくとも三つの意味で使われる。

第一に、アーキテクチャは\keyterm{分野}を指す。これは、ソフトウェア集約システムを構成するための考え方、原則、過程、方法を扱う知識領域である。建築における建築学が、建物をどう構成するかを扱うのと同じように、ソフトウェアアーキテクチャは、ソフトウェア集約システムをどう構成するかを扱う。

第二に、アーキテクチャは\keyterm{活動}を指す。アーキテクチャを作る活動は、しばしば「アーキテクチャ設計」または「アーキテクティング」と呼ばれる。ソフトウェア設計は、一般にアーキテクチャ設計、高水準設計、詳細設計に分けられる。この章は、その中でも根本的な構造や重要な意思決定に関わる部分を扱う。

第三に、アーキテクチャは\keyterm{成果}を指す。つまり、あるシステムについて、主要な構成要素、関係、原則、設計判断がどのようになっているかという結果である。この成果は、アーキテクチャ記述として表される。

かつての定義では、アーキテクチャは「システムまたは構成要素の組織構造」と理解されることが多かった。しかし、現代の考え方では、単なる構造だけでは足りない。ソフトウェアには、コード構造だけでなく、実行時構造、配置構造、情報構造、開発組織との対応、運用構造など、複数の構造がある。また、アーキテクチャは環境の中のシステムとして考える必要がある。

\begin{statementbox}{定義}
ソフトウェアアーキテクチャとは、システムの環境における根本的な概念または性質であり、構成要素、構成要素間の関係、設計と進化の原則として表れるものである。
\end{statementbox}

この定義の要点は二つである。第一に、アーキテクチャはすべての詳細を含むものではなく、根本的なものを扱う。第二に、アーキテクチャはシステムの外側も見る。利用者、組織、外部システム、ハードウェア、法律、規制、運用環境などとの関係が重要である。

\todofig{アーキテクチャの三つの意味を入れる。}{白背景、モノクロ、A4 本文幅の横長図を作成する。左から「分野としてのアーキテクチャ」「活動としてのアーキテクチャ」「成果としてのアーキテクチャ」の三つの箱を並べる。各箱に短い説明を添える。三つの箱の下に「概念・原則」「設計過程」「アーキテクチャ記述」を配置する。日本語ラベル、細線、講義ノート風。}

\subsection{利害関係者と関心事}
ソフトウェアシステムには、多くの利害関係者がいる。顧客は、いつ完成するか、いくらかかるか、投資に見合う価値があるかを気にする。利用者は、何ができるか、使いやすいか、作業を妨げないかを気にする。開発者は、構成要素の責任、インタフェース、依存関係、実装しやすさを気にする。運用担当は、配置、監視、障害対応、復旧、性能を気にする。保守担当は、変更しやすさ、影響範囲のわかりやすさ、技術的負債を気にする。

このような「気にする論点」を\keyterm{関心事}\en{concern}という。アーキテクチャの重要な役割は、利害関係者ごとの関心事を整理し、それぞれに答えられる表現を用意することである。

\begin{statementbox}{関心事の分離}
複雑なシステムを一度にすべて理解しようとすると、論点が混ざる。関心事の分離とは、性能、安全性、セキュリティ、使いやすさ、保守性などを、必要に応じて分けて考える方法である。分けることは、他の関心事を無視することではなく、順序立てて考えるための技法である。
\end{statementbox}

関心事は、機能的なもの、非機能的なもの、制約の形を取るものに分けられる。さらに、技術、事業、運用、組織、政治、経済、法令、規制、環境、社会的影響など、非常に広い範囲にまたがる。近年は、気候変動への意識の高まりにより、エネルギー効率や持続可能性もアーキテクチャ上の関心事になっている。

\begin{longtable}{P{0.25\linewidth}P{0.67\linewidth}}
\toprule
利害関係者 & 典型的な関心事 \\
\midrule
顧客・発注者 & 費用、納期、事業価値、契約適合性、長期的な運用費である。 \\
利用者 & 機能、使いやすさ、応答時間、信頼性、利用体験である。 \\
開発者 & 構成要素の責任、インタフェース、依存関係、実装容易性、テスト容易性である。 \\
運用担当 & 配置、監視、可用性、障害復旧、容量、セキュリティ運用である。 \\
保守担当 & 変更容易性、理解しやすさ、技術的負債、互換性である。 \\
監査・認証担当 & 法令・規格への適合、証拠、説明可能性、安全性、セキュリティである。 \\
\bottomrule
\end{longtable}

\todofig{利害関係者と関心事の対応図を入れる。}{白背景、モノクロ、A4 本文幅のマトリクス図を作成する。左列に「顧客」「利用者」「開発者」「運用担当」「保守担当」「監査者」を並べる。上段に「費用・納期」「機能」「性能」「可用性」「安全性」「保守性」「セキュリティ」を並べる。交点に丸印を置き、どの利害関係者がどの関心事を強く持つかを示す。日本語、講義ノート風、印刷で読みやすい。}

\subsection{アーキテクチャの用途}
アーキテクチャの第一の用途は、システムに関わる人々が共有理解を持つことである。共有理解がなければ、設計者、開発者、テスト担当、運用担当が異なる前提で作業し、結果として整合しないシステムになりやすい。

第二の用途は、設計と構築の指針を与えることである。アーキテクチャは、どの構成要素が存在するか、どの構成要素がどの責任を持つか、どのように通信するか、どの品質特性をどの方法で満たすかを示す。これにより、実装前に代替案を比較し、リスクを見つけられる。

第三の用途は、既存システムの理解である。保守、拡張、移行、改善を行うとき、既存システムのアーキテクチャを理解することが重要である。これは、逆方向の設計理解、すなわちリバースアーキテクティングに近い活動である。

また、アーキテクチャは組織構造とも関係する。コンウェイの法則は、システムの構造が、それを設計する組織のコミュニケーション構造を反映しやすいことを指摘する。組織とアーキテクチャの対応は、うまく使えば分担と責任を明確にする力になるが、悪く働くと不要な分断や複雑な依存関係を生む。

\begin{statementbox}{実務での見方}
アーキテクチャは、計画、見積り、分担、再利用、製品系列化にも使われる。単一システムの内部構造だけでなく、複数製品に共通する土台を作るためにも重要である。
\end{statementbox}

\section{ソフトウェアアーキテクチャ記述}

\subsection*{アーキテクチャ記述の役割}
\addcontentsline{toc}{subsection}{アーキテクチャ記述の役割}
小さなシステムを一人で作る場合、アーキテクチャは頭の中にあるだけでも作業できることがある。しかし、大規模で複雑なシステムをチームで作り、運用し、保守する場合には、頭の中の理解だけでは不十分である。人が入れ替わり、要求が変わり、構成が進化するため、アーキテクチャを具体的な成果物として残す必要がある。

この成果物を\keyterm{アーキテクチャ記述}\en{architecture description}という。アーキテクチャ記述は、開発者にとっては構築の設計図であり、テスト担当、品質保証担当、認証担当、運用担当、保守担当にとっては、自分の作業を進めるための根拠である。

かつては、アーキテクチャ記述は文章と非形式的な図で行われることが多かった。しかし、利害関係者と関心事が多様になるにつれて、単一の図ですべてを表すことは難しくなった。そのため、目的や読者に応じて複数のビューを使う方法が重要になった。

\subsection{アーキテクチャビューとビューポイント}
\keyterm{ビュー}\en{architecture view}とは、アーキテクチャの一面を表すものである。たとえば、機能要求をどう満たすかを示す論理ビュー、並行処理や実行時相互作用を示すプロセスビュー、システムをどこに配置するかを示す物理ビュー、実装単位や依存関係を示す開発ビューがある。

ビューを分ける理由は、関心事を分離するためである。性能を確認したい人、安全性を確認したい人、配置を考える人、コード分担を考える人は、それぞれ必要な情報が違う。一つの図にすべてを詰め込むと、誰にとっても読みにくい図になる。

\keyterm{ビューポイント}\en{architecture viewpoint}とは、ビューを作成し、解釈し、利用するための規約である。どの要素を使うか、どの関係を描くか、どの記法を使うか、どの関心事に答えるか、誰が読むかを定める。ビューポイントは、ビューの「描き方の型」である。

\begin{longtable}{P{0.25\linewidth}P{0.67\linewidth}}
\toprule
種類 & 説明 \\
\midrule
モジュールビューポイント & 実装単位としてのモジュール、その責任、依存関係、分割を表す。 \\
構成要素・接続子ビューポイント & 実行時の大きな構成要素と、それらをつなぐ通信・接続関係を表す。 \\
論理ビューポイント & 対象領域の主要概念、能力、関係を表す。 \\
シナリオ・ユースケースビューポイント & 利用者がシステムとどのように相互作用するかを表す。 \\
情報ビューポイント & 重要な情報、データの保管、アクセス、流れを表す。 \\
配置ビューポイント & ソフトウェアをどの実行環境、機器、ネットワークに置くかを表す。 \\
\bottomrule
\end{longtable}

複数のビューを使うと、別の問題が生じる。ビュー同士が矛盾する可能性である。たとえば、開発ビューでは存在する構成要素が配置ビューにはない、または論理ビューの関係が実行時ビューで説明されていない、といった問題である。これを複数ビュー問題という。対策として、ビュー間の対応関係、追跡可能性、整合性規則を明示する必要がある。

ビューを構成する方法には二つの代表的な考え方がある。\keyterm{合成的手法}では、複数のビューを別々に作り、対応規則で統合する。\keyterm{射影的手法}では、一つの統一モデルから複数のビューを取り出す。射影的手法は整合性を保ちやすいが、統一モデルが表せない関心事には弱い。

\todofig{ビューとビューポイントの関係図を入れる。}{白背景、モノクロ、A4 本文幅の概念図を作成する。中央に「アーキテクチャ記述」を置き、そこから「論理ビュー」「開発ビュー」「プロセスビュー」「配置ビュー」「シナリオビュー」へ枝分かれさせる。各ビューの上に、そのビューを定める「ビューポイント」を小さなラベルで示す。右側に「利害関係者の関心事に答える」と注記する。日本語ラベル、講義ノート風。}

\subsection{アーキテクチャパターン、様式、参照アーキテクチャ}
\keyterm{アーキテクチャ様式}\en{architectural style}とは、ソフトウェアシステムを構成する特定の作り方であり、システム全体または大きな部分に特徴を与える。たとえば、層化様式、パイプ・アンド・フィルタ、クライアントサーバ、n 層構成、マイクロサービスなどがある。

\keyterm{アーキテクチャパターン}\en{architectural pattern}とは、ソフトウェアシステムの文脈で繰り返し現れる問題に対する共通の解決法である。様式がシステム全体の大きな構成傾向を表しやすいのに対し、パターンは特定の問題への再利用可能な解決法として理解しやすい。ただし、両者の境界は厳密ではない。様式をパターンとして表すこともある。

\begin{longtable}{P{0.28\linewidth}P{0.30\linewidth}P{0.30\linewidth}}
\toprule
分類 & 例 & 初学者向けの理解 \\
\midrule
一般的な構造 & 層化、呼出し復帰、パイプ・アンド・フィルタ、黒板、サービス、マイクロサービス & 大きな構成の型である。 \\
分散システム & クライアントサーバ、n 層、ブローカ、発行購読、REST & ネットワーク越しの構成や通信の型である。 \\
方法に基づく構造 & オブジェクト指向、イベント駆動、データフロー & 設計や処理の考え方に基づく型である。 \\
利用者相互作用 & モデル・ビュー・コントローラ、表示・抽象・制御 & 画面や操作の分離に関する型である。 \\
適応システム & マイクロカーネル、反射、メタレベルアーキテクチャ & 変化や拡張に対応しやすくする型である。 \\
\bottomrule
\end{longtable}

\keyterm{参照アーキテクチャ}\en{reference architecture}とは、他のアーキテクチャを導く、または制約するアーキテクチャである。個別システム、製品系列、特定領域のシステムを作るときの共通基盤になる。自動車、医療、モノのインターネット、クラウド、航空電子、製造、通信などの領域で使われる。

\begin{statementbox}{区別の目安}
様式は「全体としてどのような形で作るか」、パターンは「繰り返す問題にどう答えるか」、参照アーキテクチャは「ある領域で個別アーキテクチャを作るための基準」である。
\end{statementbox}

\todofig{様式、パターン、参照アーキテクチャの比較図を入れる。}{白背景、モノクロ、A4 本文幅の比較表風図を作成する。列は「アーキテクチャ様式」「アーキテクチャパターン」「参照アーキテクチャ」。行は「目的」「適用範囲」「例」「実務での使い方」。各セルに短い日本語を入れる。細線、講義ノート風、印刷で読みやすい。}

\subsection{アーキテクチャ記述言語とアーキテクチャ枠組み}
\keyterm{アーキテクチャ記述言語}\en{Architecture Description Language, ADL}とは、ソフトウェアアーキテクチャを表すための領域特化言語である。ADL は、もともと大規模プログラミングにおけるモジュール相互接続言語から発展した。特定の領域や様式を対象にするものもあれば、企業全体の関心事を扱う広いものもある。

UML はソフトウェア設計で広く使われるため、アーキテクチャ記述にも使われることが多い。ADL の価値は、単に図を描くことだけではない。アーキテクチャ分析、整合性確認、場合によってはコード生成を支援できる点にある。

\keyterm{アーキテクチャ枠組み}\en{architecture framework}とは、特定領域や利害関係者共同体の中で、アーキテクチャを記述するための規約、原則、実践をまとめたものである。複数のビューポイントや ADL を組み合わせ、推奨される作業方法を定める。自動車分野の AUTOSAR、統一アーキテクチャ枠組み、オープン分散処理の参照モデルなどが例である。

\begin{statementbox}{実務での使い分け}
小規模なシステムでは、軽量な図と意思決定記録で十分なことがある。一方、複数組織、規制領域、製品系列、長期運用システムでは、ADL や枠組みを使うことで整合性と説明可能性を高められる。
\end{statementbox}

\subsection{重要な意思決定としてのアーキテクチャ}
アーキテクチャ設計は創造的な活動である。アーキテクトは、要求、品質特性、制約、利用可能な資源、組織の状況を踏まえ、多くの意思決定を行う。特に、品質特性のトレードオフは重要である。性能を高める選択が保守性を下げることがあり、セキュリティを高める選択が使いやすさや費用に影響することがある。

アーキテクチャは、意思決定の網として理解できる。ある決定が、次の決定の前提になる。したがって、重要な判断は、理由とともに記録する必要がある。

\begin{statementbox}{アーキテクチャ根拠}
アーキテクチャ根拠とは、なぜその決定をしたのかを説明する情報である。前提、検討した代替案、選択基準、トレードオフ、却下した案とその理由を含む。将来の保守担当者にとって、決定の理由は決定内容と同じくらい重要である。
\end{statementbox}

\keyterm{アーキテクチャ技術的負債}とは、現在のアーキテクチャ判断が将来の変更や運用に余分な費用を生む状態である。たとえば、初回リリースのためにモジュール性を十分考えずに作ると、後続リリースで機能追加のたびに大きな改修が必要になり、欠陥も増えやすい。この負債は、最初に判断した人ではなく、後から開発や保守を行う人が支払うことが多い。

\todofig{アーキテクチャ意思決定記録のテンプレート図を入れる。}{白背景、モノクロ、A4 本文幅の帳票風図を作成する。タイトルを「アーキテクチャ意思決定記録」とし、項目として「決定事項」「背景」「関係する要求・関心事」「検討した代替案」「選択理由」「トレードオフ」「却下した案」「将来の影響」を並べる。日本語、細線、講義ノート風。}

\section{ソフトウェアアーキテクチャ過程}

\subsection*{アーキテクチャ過程の全体像}
\addcontentsline{toc}{subsection}{アーキテクチャ過程の全体像}
アーキテクチャ設計は、単純な一方向の工程ではない。一般的には、関心事と要求を理解する分析、候補解を作る統合、候補解を検証する評価を反復する。要求が変わり、制約が明らかになり、評価で問題が見つかるたびに、分析と統合に戻る。

\subsection{文脈の中のアーキテクチャ}
アーキテクチャは真空中で作られるものではない。どの開発ライフサイクルを使うか、製品系列の中にあるか、企業全体のアーキテクチャに従うか、システム・オブ・システムズの一部かによって、作り方が変わる。

伝統的なライフサイクルでは、要求に基づいてアーキテクチャ設計段階を置くことが多い。このとき、すべての要求が同じ重要度を持つわけではない。アーキテクチャに強く影響する要求、すなわち ASR を見極める必要がある。

製品系列や製品ファミリでは、共通のアーキテクチャが先に作られ、それを土台として個別製品が作られる。個別製品は、共通部分と可変部分を持つ。共通アーキテクチャは再利用性を高めるが、過度に固定すると個別製品の自由度を下げる。

アジャイル開発では、明示的なアーキテクチャ設計段階が小さいことがあり、コードからアーキテクチャが「創発する」と説明されることもある。この考え方は、利用者中心の情報システムではうまく働く場合がある。しかし、組込みシステム、サイバーフィジカルシステム、安全性が重要なシステムでは、利用者物語に現れにくい重要品質があるため、意識的なアーキテクチャ検討が必要である。

企業アーキテクチャやシステム・オブ・システムズの文脈では、上位のアーキテクチャが個別ソフトウェアの要求や制約になる。仕様、API、準拠試験、標準インタフェースなどを通じて制約が課される。

\subsubsection{アーキテクチャと設計の関係}
アーキテクチャと設計の境界は、しばしば曖昧である。アーキテクチャはソフトウェア設計から発展してきた分野であり、設計の一部として扱われることもある。しかし、実務上は区別して考えると理解しやすい。

\begin{longtable}{P{0.22\linewidth}P{0.38\linewidth}P{0.32\linewidth}}
\toprule
観点 & アーキテクチャ & 設計 \\
\midrule
扱う範囲 & システム全体、環境、主要構成、品質特性、進化原則である。 & 構成要素の詳細、データ構造、アルゴリズム、実装方法である。 \\
要求との関係 & 要求を受けるだけでなく、利害関係者との交渉を通じて要求を形作ることがある。 & 既に定義された要求とアーキテクチャ制約を満たす方法を具体化する。 \\
変更影響 & 後から変更すると広範囲に影響しやすい。 & 比較的局所的に変更できる場合が多い。 \\
主な成果物 & アーキテクチャ記述、ビュー、意思決定記録、原則である。 & 詳細設計、インタフェース仕様、クラス設計、アルゴリズム設計である。 \\
\bottomrule
\end{longtable}

\begin{statementbox}{見分け方}
「この判断を後から変えると、多くの構成要素、品質特性、組織分担、運用方法に影響するか」と問う。影響が広ければ、アーキテクチャ判断である可能性が高い。
\end{statementbox}

\subsection{アーキテクチャ設計}
アーキテクチャ設計とは、設計原則と方法を使い、ソフトウェアアーキテクチャを作成し、記録する活動である。ここでは、システムの主要構成要素、その責任、性質、インタフェース、相互作用、環境との関係を決める。

典型的な関心事には、次のものがある。
\begin{itemize}
  \item 全体のアーキテクチャ様式と計算パラダイムである。
  \item システムを主要構成要素へ分割する方法である。
  \item 構成要素間の通信、同期、データ交換、制御の流れである。
  \item 関心事と設計責任をどの構成要素に割り当てるかである。
  \item 構成要素のインタフェースである。
  \item 性能、拡張性、資源消費、信頼性への対応である。
  \item 安全性やセキュリティのような全体横断的関心事への対応である。
\end{itemize}

アーキテクチャ設計は、\keyterm{分析}、\keyterm{統合}、\keyterm{評価}の三つの活動を反復する。

\subsubsection{アーキテクチャ分析}
アーキテクチャ分析は、アーキテクチャ上重要な要求、すなわち ASR を集め、整理する活動である。ASR は、アーキテクチャに影響する要求である。たとえば、厳しい応答時間、大量アクセスへの拡張性、高い可用性、法規制への適合、強いセキュリティ、外部システムとの接続方式などが該当する。

分析では、既知の要求、利害関係者のニーズ、環境制約、上位アーキテクチャ、利用可能な資源を確認する。要求と制約の組合せが実現困難な場合には、利害関係者と交渉し、期待や要求を調整する必要がある。

\subsubsection{アーキテクチャ統合}
アーキテクチャ統合は、分析で見つかった問題に対して候補解を作る活動である。構成要素をどう分けるか、どの様式を採るか、どの技術を使うか、どの品質特性をどの戦術で満たすかを検討する。

ここで重要なのは、局所的に最適な解が全体として最適とは限らないことである。性能を高める候補解が変更容易性を下げることがある。セキュリティを高める候補解が使いやすさを下げることがある。統合は、複数の関心事の相互作用を考慮する活動である。

\subsubsection{アーキテクチャ評価}
アーキテクチャ評価は、選んだ候補解が ASR を満たすかを確認し、手戻りが必要な箇所を見つける活動である。評価は最後に一度だけ行うものではない。分析と統合の途中で繰り返し行い、問題を早く見つけるほど、修正費用を抑えられる。

\todofig{アーキテクチャ分析・統合・評価の反復サイクルを入れる。}{白背景、モノクロ、A4 本文幅の循環図を作成する。中央に「アーキテクチャ設計」を置き、周囲に「分析：ASR と制約を整理」「統合：候補解を作る」「評価：満たすか確認」を円形に配置する。外側から「利害関係者の関心事」「要求」「環境制約」が分析へ入る矢印を描く。評価から分析へ戻る反復矢印を描く。日本語ラベル、講義ノート風。}

\subsection{アーキテクチャの実践、方法、戦術}
アーキテクチャの実践には、多くの方法がある。方法とは、関心事の発見、要求整理、候補解作成、文書化、レビュー、評価をどの順序で行うかを定める作業手順である。戦術とは、特定の品質特性を満たすための小さな設計上の手段である。

たとえば、可用性を高めるには冗長化、監視、障害切替が使われることがある。性能を高めるにはキャッシュ、非同期処理、並列化が使われることがある。変更容易性を高めるには抽象化、疎結合、情報隠蔽が使われることがある。セキュリティを高めるには認証、認可、入力検証、監査ログが使われることがある。

\begin{statementbox}{注意}
戦術は魔法の解決策ではない。どの戦術にも副作用がある。キャッシュは性能を高めるが、整合性管理を難しくすることがある。冗長化は可用性を高めるが、費用と運用複雑性を増やすことがある。
\end{statementbox}

\subsection{大規模なアーキテクチャ活動}
アーキテクチャ設計は、ライフサイクル上の一段階を指すことがある。しかし、広い意味のアーキテクティングは、それだけではない。大規模な組織やシステムでは、個別システムのアーキテクチャが、企業アーキテクチャ、製品系列アーキテクチャ、システム・オブ・システムズのアーキテクチャと関係する。

このような文脈では、次の活動が重要になる。
\begin{itemize}
  \item \keyterm{アーキテクチャ実装管理}：実装がアーキテクチャに従っているかを確認する。
  \item \keyterm{アーキテクチャ保守}：実装後もアーキテクチャを管理し、拡張する。
  \item \keyterm{アーキテクチャ管理}：組織内の複数アーキテクチャの関係を管理する。
  \item \keyterm{アーキテクチャ知識管理}：意思決定、教訓、仕様、文書などを再利用可能な資産として残す。
\end{itemize}

大規模なアーキテクチャ活動では、個別システムの最適化だけでは不十分である。組織全体で一貫性、相互運用性、再利用性、進化可能性を保つ必要がある。

\section{ソフトウェアアーキテクチャ評価}

\subsection{アーキテクチャにおける「よさ」}
アーキテクチャの「よさ」は、一つの尺度だけでは測れない。ある利害関係者にとってよいアーキテクチャが、別の利害関係者にとっては問題を含むことがある。たとえば、セキュリティを非常に強くした構成は安全に見えるが、費用が高く、検証が難しく、利用者の操作が複雑になることがある。

建築の古典的な考え方では、建物には強さ、用、見た目のよさが必要だとされる。この考え方をソフトウェアに置き換えると、次の問いになる。
\begin{itemize}
  \item 長期運用と進化に耐えられるか。
  \item 意図した用途に合っているか。
  \item このアーキテクチャで、費用対効果よく構築できるか。
  \item 構築、利用、保守を行う人にとって明確で理解しやすいか。
\end{itemize}

\keyterm{アーキテクチャトレードオフ分析法}\en{Architecture Tradeoff Analysis Method, ATAM}は、品質特性に基づいてアーキテクチャを評価する方法である。品質特性を整理し、シナリオを使って評価し、品質要求同士のトレードオフを分析する。ほかにも、ソフトウェアアーキテクチャ分析法\en{SAAM}や品質特性ワークショップ\en{QAW}などの方法がある。

\begin{statementbox}{評価の本質}
評価の目的は、アーキテクチャに点数を付けることだけではない。重大なリスク、隠れた前提、品質特性間の衝突、説明できない判断を早く発見することである。
\end{statementbox}

\subsection{アーキテクチャについての推論}
アーキテクチャ評価は、明確なアーキテクチャ記述があるほど効果的である。ビュー、モデル、意思決定記録があれば、それらを調べ、質問し、分析できる。信頼性、安全性、セキュリティのような専門的関心事では、それぞれの分野で使われる専門的な表現が必要になることがある。

しかし、現実にはアーキテクチャ文書が未完成、不完全、古い、存在しない場合も多い。その場合、評価は関係者の知識に依存する。これは危険である。なぜなら、人の記憶は抜け落ちやすく、関係者が退職すると知識も失われるからである。

ユースケースは、アーキテクチャの完全性と整合性を確認する手段として使える。ユースケースの各ステップを見て、そのステップを実現する構成要素がアーキテクチャに存在するか、必要な通信やデータが説明されているかを確認する。

\todofig{ユースケースからアーキテクチャ要素をたどる図を入れる。}{白背景、モノクロ、A4 本文幅の横長図を作成する。左に「ユースケースの手順」を縦に並べ、右に「構成要素」「インタフェース」「データ」「配置先」を並べる。各手順から関連するアーキテクチャ要素へ矢印を引く。下に「対応しない手順は不足の候補」「使われない要素は過剰の候補」と注記する。日本語ラベル、講義ノート風。}

\subsection{アーキテクチャレビュー}
アーキテクチャレビューは、アーキテクチャの状態、品質、リスクを確認するための有効な方法である。レビューは、非公式な専門家レビューとして行われることもあれば、チェックリストに基づいて構造化されることもある。

より能動的なレビューでは、単に「この項目を確認したか」と聞くのではなく、レビュアが特定の活動を行って情報を得る。たとえば、「障害が発生したときに、どの構成要素がどの順で影響を受けるかを説明する」「利用者数が 10 倍になったときに、どの構成要素がボトルネックになるかを示す」といった活動である。

\begin{longtable}{P{0.25\linewidth}P{0.67\linewidth}}
\toprule
レビュー観点 & 確認すること \\
\midrule
要求対応 & ASR に答える構成や判断があるか。 \\
品質特性 & 性能、可用性、変更容易性、安全性、セキュリティなどが説明されているか。 \\
トレードオフ & ある品質を優先したことで、別の品質にどのような影響があるか。 \\
リスク & 未検証の前提、難しい技術、外部依存、単一障害点があるか。 \\
説明可能性 & 重要な意思決定に根拠が残っているか。 \\
実装適合性 & 実装がアーキテクチャに従っているか、または逸脱が管理されているか。 \\
\bottomrule
\end{longtable}

\subsection{アーキテクチャ測度}
\keyterm{アーキテクチャ測度}\en{architecture metrics}とは、アーキテクチャの性質を数量で表すものである。測度は万能ではないが、複雑さ、依存関係、変更リスク、規約適合性を把握する手がかりになる。

代表的な測度には、構成要素間依存、循環依存、循環的複雑度、内部モジュール複雑度、結合度、凝集度、ネストの深さ、パターン・様式・API への準拠などがある。

継続的開発や DevOps の文脈では、アーキテクチャそのものではなく、アーキテクチャの状態を反映する過程測度も使われる。変更リードタイム、配置頻度、サービス復旧平均時間、変更失敗率などである。変更が遅く、復旧が難しく、失敗率が高い場合、アーキテクチャが変更や運用に適していない可能性がある。

\todofig{アーキテクチャ評価の観点マップを入れる。}{白背景、モノクロ、A4 本文幅の放射状図を作成する。中央に「アーキテクチャ評価」を置き、周囲に「よさ」「推論」「レビュー」「測度」「トレードオフ」「リスク」を配置する。各項目に短い説明を付ける。日本語ラベル、細線、講義ノート風。}

\section{トピックと参考文献の対応}
この表は、SWEBOK が示す「トピック対参考資料の行列」を、学習用に簡略化したものである。詳しく学ぶときは、各トピックに対応する文献を読む。

\begin{longtable}{P{0.42\linewidth}P{0.50\linewidth}}
\toprule
トピック & 主な参考資料 \\
\midrule
1. ソフトウェアアーキテクチャの基礎 & Bass et al.; Maier and Rechtin; Rozanski and Woods; Taylor et al. \\
1.1 「アーキテクチャ」という語の意味 & Bass et al.; Budgen; Maier and Rechtin \\
1.2 利害関係者と関心事 & Bass et al.; Rozanski and Woods; Taylor et al. \\
1.3 アーキテクチャの用途 & Bass et al.; Rozanski and Woods \\
2. アーキテクチャ記述 & Bass et al.; Rozanski and Woods; Sommerville; Taylor et al. \\
2.1 ビューとビューポイント & Budgen; Maier and Rechtin; Rozanski and Woods; Sommerville \\
2.2 パターン、様式、参照アーキテクチャ & Bass et al.; Budgen; Rozanski and Woods; Sommerville; Taylor et al. \\
2.3 アーキテクチャ記述言語と枠組み & Bass et al.; Maier and Rechtin; Rozanski and Woods; Taylor et al. \\
2.4 重要な意思決定としてのアーキテクチャ & Rozanski and Woods; Sommerville \\
3. アーキテクチャ過程 & Maier and Rechtin; Rozanski and Woods; Taylor et al. \\
3.1 文脈の中のアーキテクチャ & Maier and Rechtin; Taylor et al. \\
3.1.1 アーキテクチャと設計の関係 & Sommerville; Taylor et al. \\
3.2 アーキテクチャ設計 & Bass et al. \\
3.2.1 アーキテクチャ分析 & Bass et al.; Taylor et al. \\
3.2.2 アーキテクチャ統合 & Bass et al. \\
3.2.3 アーキテクチャ評価 & Bass et al.; Rozanski and Woods \\
3.3 実践、方法、戦術 & Bass et al.; Maier and Rechtin; Rozanski and Woods \\
3.4 大規模なアーキテクチャ活動 & Maier and Rechtin; Sommerville \\
4. アーキテクチャ評価 & Bass et al.; Rozanski and Woods; Taylor et al. \\
4.1 アーキテクチャの「よさ」 & Bass et al.; Budgen \\
4.2 アーキテクチャについての推論 & Rozanski and Woods \\
4.3 アーキテクチャレビュー & Bass et al. \\
4.4 アーキテクチャ測度 & Bass et al. \\
\bottomrule
\end{longtable}

\section{発展的な読み物}
SWEBOK 02 章で紹介される発展的な読み物を、日本語で読む目的に合わせて整理する。

\begin{longtable}{P{0.32\linewidth}P{0.60\linewidth}}
\toprule
文献 & 読むと得られること \\
\midrule
Perry and Wolf & ソフトウェアアーキテクチャ研究の基礎的な考え方を学べる。複数ビュー、アーキテクチャ様式、設計との違いなどの出発点を理解できる。 \\
Bass, Clements and Kazman & 品質特性、アーキテクチャ設計、記述、評価、技術的負債を実務的に学べる。 \\
Kruchten & 4+1 ビューモデルを学べる。論理、開発、プロセス、物理の四つのビューを、シナリオで結びつける考え方である。 \\
Rozanski and Woods & 利害関係者、関心事、ビュー、ビューポイントを使ってアーキテクチャを記述する実践的な方法を学べる。 \\
Taylor, Medvidović and Dashofy & アーキテクチャの基礎、理論、実践を包括的に学べる。実装、配置、非機能特性、信頼とセキュリティも扱う。 \\
Clements et al. & ビューとその種類を使って、アーキテクチャをどう文書化するかを学べる。 \\
Brown & 開発者の視点から、アーキテクチャドライバ、品質、制約、リスク管理を学べる。 \\
Fairbanks & リスク駆動で「必要十分な」アーキテクチャを行う考え方を学べる。軽量でアジャイルな文脈にも合う。 \\
Erder, Pureur and Woods & アジャイル、クラウド、DevOps 時代の継続的アーキテクチャを学べる。セキュリティ、回復性、拡張性にも重点がある。 \\
\bottomrule
\end{longtable}

\section{参考文献}
{\small\sloppy
\begin{enumerate}[label={[\arabic*]},leftmargin=2.5em]
\item M. Ali Babar and I. Gorton, ``Software Architecture Review: The State of the Practice,'' IEEE Computer, July 2009.
\item L. Bass, P. Clements, and R. Kazman, \textit{Software Architecture in Practice}, 4th edition, 2021.
\item L. Bass, J. Ivers, M. H. Klein, and P. Merson, \textit{Reasoning Frameworks}, CMU/SEI-2005-TR-007, 2005.
\item F. Brooks, \textit{The Design of Design}, Addison-Wesley, 2010.
\item S. Brown, \textit{Software Architecture for Developers}, Leanpub, 2018.
\item D. Budgen, \textit{Software Design: Creating Solutions for Ill-Structured Problems}, 3rd edition, CRC Press, 2021.
\item F. Buschmann, R. Meunier, H. Rohnert, P. Sommerlad, and M. Stal, \textit{Pattern-Oriented Software Architecture}, John Wiley \& Sons, 1996.
\item H. Cervantes and R. Kazman, \textit{Designing Software Architectures: A Practical Approach}, 2nd edition, Addison-Wesley, 2024.
\item P. Clements et al., \textit{Documenting Software Architecture: Views and Beyond}, 2nd edition, Addison-Wesley, 2011.
\item P. Clements, R. Kazman, and M. Klein, \textit{Evaluating Software Architectures}, Addison-Wesley, 2001.
\item M. E. Conway, ``How Do Committees Invent?'' \textit{Datamation}, 14(4), 28--31, 1968.
\item E. W. Dijkstra, ``On the role of scientific thought,'' 1974.
\item T. Erl, \textit{SOA Design Patterns}, Prentice-Hall, 2009.
\item P. Eeles and P. Cripps, \textit{The Process of Software Architecting}, Addison-Wesley, 2010.
\item M. Erder, P. Pureur, and E. Woods, \textit{Continuous Architecture in Practice: Software Architecture in the Age of Agility and DevOps}, Addison-Wesley, 2021.
\item G. Fairbanks, \textit{Just Enough Software Architecture: A Risk-Driven Approach}, Marshall \& Brainerd, 2010.
\item E. Fernandez-Buglioni, \textit{Security Patterns in Practice: Designing Secure Architectures Using Software Patterns}, Wiley, 2013.
\item R. T. Fielding and R. N. Taylor, ``Principled design of the modern web architecture,'' \textit{ACM Transactions on Internet Technology}, 2(2), 115--150, 2002.
\item M. Fowler, D. Rice, M. Foemmel, E. Hieatt, R. Mee, and R. Stafford, \textit{Patterns of Enterprise Application Architecture}, Addison-Wesley, 2003.
\item C. Hofmeister, P. B. Kruchten, R. L. Nord, H. Obbink, A. Ran, and P. America, ``A general model of software architecture design derived from five industrial approaches,'' \textit{The Journal of Systems and Software}, 80, 106--126, 2007.
\item C. Hofmeister, R. L. Nord, and D. Soni, \textit{Applied Software Architecture}, Addison-Wesley, 2000.
\item ISO/IEC/IEEE 24765:2017, \textit{Systems and Software Engineering - Vocabulary}, 2nd edition, 2017.
\item ISO/IEC/IEEE 42010:2011, \textit{Systems and software engineering - Architecture description}, 2011.
\item R. Kazman, S. Haziyev, A. Yakuba, and D. A. Tamburri, ``Managing Energy Consumption as an Architectural Quality Attribute,'' \textit{IEEE Software}, 35(5), 102--107, 2018.
\item P. B. Kruchten, ``The 4+1 View Model of Architecture,'' \textit{IEEE Software}, 12(6), 1995.
\item P. B. Kruchten, R. L. Nord, and I. Ozkaya, \textit{Managing Technical Debt: Reducing Friction in Software Development}, Addison-Wesley, 2019.
\item Z. Li, P. Liang, and P. Avgeriou, ``Architecture viewpoints for documenting architectural technical debt,'' \textit{Software Quality Assurance}, Elsevier, 2016.
\item A. MacCormack, J. Rusnak, and C. Baldwin, ``Exploring the Duality between Product and Organizational Architectures: A Test of the `Mirroring' Hypothesis,'' \textit{Research Policy}, 41, 1309--1324, 2012.
\item M. W. Maier and E. Rechtin, \textit{The Art of Systems Architecting}, 3rd edition, CRC Press, 2021.
\item N. Medvidović, D. S. Rosenblum, D. F. Redmiles, and J. E. Robbins, ``Modeling software architectures in the Unified Modeling Language,'' \textit{ACM Transactions on Software Engineering and Methodology}, 11(1), 2--57, 2002.
\item H. Obbink et al., \textit{Report on Software Architecture Review and Assessment (SARA)}, version 1.0, 2002.
\item D. L. Parnas, ``On the criteria to be used in decomposing systems into modules,'' \textit{Communications of the ACM}, 15(12), 1053--1058, 1972.
\item D. L. Parnas and D. M. Weiss, ``Active Design Reviews: Principles and Practices,'' Proceedings of the 8th International Conference on Software Engineering, 215--222, 1985.
\item D. Perry and A. Wolf, ``Foundations for the study of software architecture,'' \textit{ACM SIGSOFT Software Engineering Notes}, 17(4), 40--52, 1992.
\item E. Poort and H. van Vliet, ``RCDA: Architecting as a Risk- and Cost Management Discipline,'' \textit{Journal of Systems and Software}, 2012.
\item R. Prieto-Diaz and J. M. Neighbors, ``Module Interconnection Languages,'' \textit{Journal of Systems and Software}, 6(4), 307--334, 1986.
\item C. Richardson, \textit{Microservices Patterns}, Manning Publications, 2019.
\item N. Rozanski and E. Woods, \textit{Software Systems Architecture: Working with Stakeholders Using Viewpoints and Perspectives}, 2nd edition, Addison-Wesley, 2011.
\item M. Shaw and D. Garlan, \textit{Software Architecture: Perspectives on an Emerging Discipline}, Prentice Hall, 1996.
\item I. Sommerville, \textit{Software Engineering}, 10th edition, 2016.
\item R. N. Taylor, N. Medvidović, and E. Dashofy, \textit{Software Architecture: Foundations, Theory, and Practice}, Wiley, 2009.
\item R. Weinreich and G. Buchgeher, ``Towards supporting the software architecture life cycle,'' \textit{The Journal of Systems and Software}, 85, 546--561, 2012.
\end{enumerate}
}

\section{画像 TODO 一覧}
本文に配置した画像 TODO を再掲する。実際に図を作る場合は、各 TODO のプロンプトをそのまま画像生成に使うと、A4 資料に合う図を作りやすい。
\begin{enumerate}
  \item 02 章の全体地図。
  \item アーキテクチャの三つの意味。
  \item 利害関係者と関心事の対応図。
  \item ビューとビューポイントの関係図。
  \item 様式、パターン、参照アーキテクチャの比較図。
  \item アーキテクチャ意思決定記録のテンプレート図。
  \item アーキテクチャ分析・統合・評価の反復サイクル。
  \item ユースケースからアーキテクチャ要素をたどる図。
  \item アーキテクチャ評価の観点マップ。
\end{enumerate}

\section{章末確認問題}
\begin{enumerate}
  \item 「アーキテクチャ」という語の三つの意味を説明せよ。
  \item 利害関係者と関心事の違いを、自分の言葉で説明せよ。
  \item なぜアーキテクチャ記述では複数のビューを使うのか。
  \item ビューとビューポイントの違いを説明せよ。
  \item アーキテクチャ様式、アーキテクチャパターン、参照アーキテクチャを区別せよ。
  \item アーキテクチャ上重要な要求 ASR の例を三つ挙げよ。
  \item アーキテクチャ根拠を記録する理由を説明せよ。
  \item アーキテクチャ評価で、トレードオフを確認することが重要な理由を説明せよ。
  \item アーキテクチャ測度の例を三つ挙げ、それぞれ何を知るためのものか説明せよ。
\end{enumerate}

\begin{statementbox}{解答の方向性}
1 は三義、2 は主体と論点、3 と 4 は関心事別の表現、5 は型・解決法・領域基準の違いで説明する。6 以降は、性能、可用性、セキュリティ、変更容易性、法規制、外部連携を例にできる。
\end{statementbox}

\end{document}
